\documentclass[10pt,twocolumn,letterpaper]{article}

\usepackage{iccv}
\usepackage{times}
\usepackage{epsfig}
\usepackage{graphicx}
\usepackage{amsmath}
\usepackage{amssymb}

% Include other packages here, before hyperref.

% If you comment hyperref and then uncomment it, you should delete
% egpaper.aux before re-running latex.  (Or just hit 'q' on the first latex
% run, let it finish, and you should be clear).
\usepackage[breaklinks=true,bookmarks=false]{hyperref}

\iccvfinalcopy % *** Uncomment this line for the final submission

\def\iccvPaperID{****} % *** Enter the ICCV Paper ID here
\def\httilde{\mbox{\tt\raisebox{-.5ex}{\symbol{126}}}}

% Pages are numbered in submission mode, and unnumbered in camera-ready
%\ificcvfinal\pagestyle{empty}\fi
\setcounter{page}{1}
\begin{document}

%%%%%%%%% TITLE
\title{SPARE: SAM3 Prompt Analysis and Robustness Evaluation}

\author{Yara Hisham \and Aliaa Gheis \and Nouran Hany \and Sarah Elzayat \and Maryam Muhammed 
\\
\\
Computer Engineering Department\\
Cairo University Faculty of Engineering, Cairo, Egypt
}

\maketitle
%\thispagestyle{empty}


%%%%%%%%% ABSTRACT
\begin{abstract}
The rise of Segment Anything, SAM~\cite{kirillov2023sam}, with its versions has greatly impacted the image segmentation industry. 
Especially SAM~3~\cite{peng2025sam3}, whose ability to have both geometric and text prompts made a significant impact. 
However, there is no systematic study that analyzes SAM~3~\cite{peng2025sam3} performance on degraded visuals or text prompts. 
We introduce SPARE, the first stress-testing framework that benchmarks SAM~3~\cite{peng2025sam3} under both visual degradations and typographic corruptions. 
Our analysis has shown that, despite slight performance degradation under visual variations, SAM~3~\cite{peng2025sam3} is highly sensitive and prone to very small changes in text prompts.
\end{abstract}

%%%%%%%%% BODY TEXT
\section{Introduction}

Image segmentation has long been treated as a purely visual task, where model inputs are 
clean images and prompts are precise spatial coordinates. The release of SAM~\cite{kirillov2023sam} 
and its successors challenged this assumption by enabling flexible, zero-shot segmentation 
through points, boxes, and masks. SAM~3~\cite{peng2025sam3} takes this further, accepting 
natural language concept prompts alongside spatial cues—a shift that makes the model 
significantly more expressive, but also more exposed: it must now interpret potentially 
noisy or ambiguous text in addition to potentially degraded images.

Despite this expanded attack surface, no systematic study has examined how SAM~3 behaves 
under realistic input noise. Prior robustness work either targets earlier SAM 
versions~\cite{wang2023empirical,chen2024robustsam}, focuses on spatial prompt 
jitter~\cite{moskalenko2026breps}, or studies linguistic sensitivity in classification and 
retrieval models~\cite{materzynska2025scam,schmalfuss2025parc}—none of which address the 
unique failure modes introduced by concept-level prompting in a segmentation model.

This paper presents \textbf{SPARE}, a stress-testing framework that fills this gap. 
We make four contributions:
\begin{enumerate}
    \item We establish a dual-pipeline evaluation protocol that measures output consistency 
    rather than absolute accuracy, removing dependence on fixed human annotations.
    \item We benchmark SAM~3 under visual degradations including Gaussian noise and 
    horizontal flips across multiple severity levels.
    \item We evaluate SAM~3's concept-prompt pipeline under realistic typographic 
    corruptions—character substitutions, deletions, and naturalistic misspellings.
    \item We provide a direct, controlled comparison between visual and linguistic 
    fragility within the same model, revealing where multimodal grounding introduces 
    the greater operational risk.
\end{enumerate}

The paper outline is as follows: Section~\ref{sec:related} reviews 
related work. Section~\ref{sec:method} describes the SPARE framework and corruption 
protocols. Section~\ref{sec:expSetup} states the setup we used to produce the results.  Section~\ref{sec:results} presents experimental results, and 
Section~\ref{sec:conc} concludes with a short discussion of findings, limitations, and 
directions for future work.

%------------------------------------------------------------------------
\section{Related Work}
\label{sec:related}

\subsection{Promptable Segmentation}

The Segment Anything Model (SAM)~\cite{kirillov2023sam} established a new paradigm for instance segmentation by supporting multiple kinds of prompts, which include: points, bounding boxes, and coarse masks without task-specific fine-tuning. SAM~2~\cite{ravi2024sam2} extended this framework to video through memory-augmented temporal propagation, preserving spatial masks across frames. Most recently, SAM~3~\cite{peng2025sam3} introduced explicit concept-level natural language prompting alongside traditional spatial cues, effectively bridging low-level localization with high-level semantic grounding and evaluated on the SACo-Gold benchmark~\cite{meta2025saco}. While this progression reflects growing architectural sophistication and flexibility, it also introduces qualitatively new failure modes: the model must now reconcile noisy or ambiguous \emph{linguistic} inputs with visual evidence a vulnerability that purely spatial segmenters do not face.

\subsection{Robustness and Corruption Benchmarks}

Systematic evaluation of neural networks under distribution shift and synthetic corruption has become a cornerstone of reliability assessment. ImageNet-C~\cite{hendrycks2019benchmarking} established the canonical protocol of applying pixel-level distortions such as Gaussian noise, blur families, and environmental artifacts at graduated severity levels to measure degradation curves. This methodology has since been adapted to web-scale multi-modal settings; LAION-C~\cite{li2025laionc}, for instance, applies analogous corruptions to evaluate generalized vision-language encoders.

Within the segmentation domain specifically, Wang et al.~\cite{wang2023empirical} conducted early audits of SAM's sensitivity to low-quality inputs, revealing pronounced vulnerabilities that motivated RobustSAM~\cite{chen2024robustsam}, a drop-in architectural extension designed to recover performance on degraded imagery. Separately, the BREPS framework~\cite{moskalenko2026breps} formalized how bounding-box prompt inaccuracies translate into segmentation quality degradation. A critical limitation unifies these efforts: they evaluate absolute mask accuracy against static human annotations under a single noise modality, leaving the \emph{output consistency} of concept-driven models under paired perturbations unexamined.

\subsection{Linguistic Sensitivity and Prompt Robustness in VLMs}

As vision models increasingly depend on textual guidance, they inherit the well-documented prompt sensitivity of vision-language models (VLMs). Systematic studies have shown that subtle changes in phrasing, formatting, or surface orthography can meaningfully alter model behavior. The SCAM benchmark~\cite{materzynska2025scam} demonstrated that typographic noise—ranging from systematic character substitutions to naturalistic misspellings—degrades the reasoning chains of multimodal foundation models far more severely than clean-text baselines would predict. Complementarily, PARC~\cite{schmalfuss2025parc} exposed systematic geometric asymmetries in how VLMs resolve referring expressions under structural linguistic variation. Together, these works establish that linguistic degradation is a first-class robustness concern, not merely a pre-processing artifact—yet neither framework considers models whose primary output is a spatial segmentation mask rather than a classification or retrieval result.

\subsection{Positioning SPARE}
Prior work falls into three separate areas: benchmarking visual corruptions~\cite{chen2024robustsam,li2025laionc}, 
studying spatial prompt noise~\cite{moskalenko2026breps}, and probing linguistic sensitivity in 
VLMs~\cite{materzynska2025scam,schmalfuss2025parc}. SPARE brings these threads together in a single framework. Instead of measuring how much accuracy drops relative to human annotations, we compare model outputs before and after perturbation, if the same input produces consistent results under noise, the model is robust. This lets us directly ask: does SAM~3 break more easily when its 
\emph{image} is degraded, or when its \emph{text prompt} is corrupted?

%-------------------------------------------------------------------------
% \subsection{Illustrations, graphs, and photographs}

% All graphics should be centered.  Please ensure that any point you wish to
% make is resolvable in a printed copy of the paper.  Resize fonts in figures
% to match the font in the body text, and choose line widths which render
% effectively in print.  Many readers (and reviewers), even of an electronic
% copy, will choose to print your paper in order to read it.  You cannot
% insist that they do otherwise, and therefore must not assume that they can
% zoom in to see tiny details on a graphic.

% When placing figures in \LaTeX, it's almost always best to use
% \verb+\includegraphics+, and to specify the  figure width as a multiple of
% the line width as in the example below
% {\small\begin{verbatim}
%    \usepackage[dvips]{graphicx} ...
%    \includegraphics[width=0.8\linewidth]
%                    {myfile.eps}
% \end{verbatim}
% }
% \begin{table}
% \begin{center}
% \begin{tabular}{|l|c|}
% \hline
% Method & Frobnability \\
% \hline\hline
% Theirs & Frumpy \\
% Yours & Frobbly \\
% Ours & Makes one's heart Frob\\
% \hline
% \end{tabular}
% \end{center}
% \caption{Results.   Ours is better.}
% \end{table}
%------------------------------------------------------------------------
\section{Methodology}
\label{sec:method}
SPARE evaluates SAM~3~\cite{peng2025sam3} under systematic input corruption across two complementary pipelines that exercise different prompt modalities. Rather than comparing model outputs against ground-truth annotations,we adopt an \emph{output-stability} paradigm: we measure how consistently SAM~3 produces the same segmentation when inputs are progressively
corrupted, relative to its own clean-input predictions.

We apply three degradation families: \emph{linguistic noise} on text prompts,
\emph{Gaussian image noise}, and \emph{spatial flips}.  A fourth family---geometric prompt jitter---was implemented but deferred; see
\S\ref{sec:geo-deferred}.


%--------------------------------------------------------------------

\subsection{Evaluation Pipelines}
%SA-1B
The first pipeline is invoking SAM~3 with a fixed generic text prompt (\texttt{"visual"}), causing it to operate in general instance-segmentation mode. Images are drawn from
SA-1B~\cite{kirillov2023sam}. Only image-level degradations apply to this pipeline (Gaussian noise and spatial flips); the prompt is constant and cannot
be meaningfully corrupted.
%SACo-Gold

The second pipeline is invoking SAM~3 with the concept-level text prompt stored in each SACo-Gold record (e.g.\ \texttt{"wheel"}, \texttt{"tree branch"}), drawn from the SA-1B captioner subset of SACo-Gold~\cite{meta2025saco} (13,258 image concept
pairs). All three degradation families are applied to this pipeline.
%--------------------------------------------------------------------
\subsection{Prompt Selection for the SA-1B Pipeline}
\label{sec:prompt-selection}

The SA-1B pipeline invokes SAM~3 with a single fixed generic text prompt
intended to trigger \emph{general} instance-segmentation mode rather than concept-specific retrieval.  To determine the most appropriate prompt, we evaluated three candidates---\texttt{"visual"}, \texttt{"things"}, and \texttt{"objects"} on the full 10,548 image SA-1B subset by running a single forward pass per image with each prompt and measuring (i)~prediction coverage (fraction of images that receive at least one predicted mask) and (ii)~mean IoU against SA-1B ground-truth annotations.

Results are summarised in Table~\ref{tab:prompt-selection}.
\texttt{"things"} and \texttt{"objects"} caused SAM~3 to return \emph{no
predictions} for 99.97\% and 96.8\% of images respectively, making them
unsuitable as a baseline.  \texttt{"visual"} produced predictions on
89.5\% of images (9,443 of 10,548) with a mean mIoU of 0.737 against
ground truth.  Accordingly, \texttt{"visual"} was adopted as the fixed
prompt for all SA-1B pipeline experiments.

\begin{table}[h]
\centering
\small
\begin{tabular}{lccc}
\hline
Prompt & \#~Evaluated & Skipped (no pred.) & mIoU \\
\hline
\texttt{visual}  & 9{,}443 &  1{,}105 & 0.737 \\
\texttt{things}  &       3 & 10{,}545 & 0.720 \\
\texttt{objects} &     335 & 10{,}213 & 0.776 \\
\hline
\end{tabular}
\vspace{3 pt}
\caption{Prompt selection results on the SA-1B subset.
         Skipped entries are images for which SAM~3 returned no predicted
         masks under the given prompt.  \texttt{visual} is selected for
         its near-universal coverage.}
\label{tab:prompt-selection}
\end{table}
%---------------------------------------------------------------------
\subsection{Datasets}
\label{sec:datasets}
  \textbf{SA-1B}~\cite{kirillov2023sam} is the large-scale segmentation dataset released  alongside SAM, comprising 11M images and 1.1B high-quality masks. We use it as the source for the SA-1B pipeline, where SAM~3 is run in general segmentation mode with a fixed generic prompt. Our evaluation subset consists of 10,548 images at baseline; after matching clean and degraded predictions,
  9,442--9,443 valid pairs remain per degradation level (1,105 entries are skipped due to absent baseline matches).


  \textbf{SACo-Gold}~\cite{meta2025saco} is the official SAM~3 evaluation benchmark,containing 156K image--concept pairs triple-annotated across seven subsets. We use the \emph{SA-1B captioner} subset, which provides 13,258 image concept pairs and 30,306 masks derived from natural-language captions of SA-1B images.
  This subset drives the SACo-Gold pipeline, where SAM~3 is explicitly prompted with the stored concept label. After matching, 8,046--8,048 valid pairs remain
  per degradation level (5,210 entries skipped due to absent baseline matches).

%--------------------------------------------------------------------
\subsection{Linguistic Degradation}
Given a prompt $\mathbf{p}$, exactly $\ell$ non-space character positions are
selected uniformly at random (without replacement) and each character $c$ is replaced with a uniformly sampled lowercase letter $c' \neq c$.  Spaces are
protected so word boundaries are preserved.  We test three levels $\ell \in \{1, 2, 3\}$, corresponding to a single-key slip, a moderately misspelled word, and a heavily corrupted label. The source image is held fixed. This degradation is applied \emph{exclusively} to the SACo-Gold pipeline.

% \paragraph{Prompt-length stratification.}
Because $\ell$ is an absolute character count, its \emph{effective} corruption intensity depends on prompt length.  A single substitution on a 4-character prompt corrupts 25\% of its characters, whereas the
same substitution on a 15-character phrase corrupts only 7\%.  To test whether SAM~3's linguistic fragility reflects this asymmetry, we partition SACo-Gold prompts into three word-count buckets: \textbf{1-word} (1{,}451 pairs, mean
5.8 non-space characters), \textbf{2-word} (4{,}665 pairs, mean 9.1), and
\textbf{3\kern0pt+\kern0pt-word} (7{,}142 pairs, mean 17.0).  For each bucket
and each level $\ell$, we compute mIoU and SI independently, and report the
\emph{effective corruption rate}
%
\begin{equation}
  \eta(\ell, b) = \frac{\ell}{\bar{c}_b},
  \label{eq:eff-corruption}
\end{equation}
%
where $\bar{c}_b$ is the mean non-space character count of bucket $b$.
$\eta$ normalizes the noise level by prompt length, making fragility
comparable between buckets on the same footing.  Results appear in
\S\ref{sec:results-text-length}.
%--------------------------------------------------------------------
\subsection{Visual Degradation: Gaussian Noise}
The source image is corrupted channel-wise as:
%
\begin{equation}
  \mathbf{x}' = \mathrm{clip}(\mathbf{x} + \boldsymbol{\varepsilon},\;0,\;255),
  \quad
  \boldsymbol{\varepsilon} \sim \mathcal{N}(\mathbf{0},\,\sigma^2\mathbf{I}).
  \label{eq:awgn}
\end{equation}
%
We test $\sigma \in \{10, 25, 50\}$, representing mild, moderate, and severe
sensor degradation.  All prompts are held fixed.  This degradation is applied to \emph{both} pipelines.

%Cross-modal comparison.


Applying identical Gaussian corruptions to both pipelines lets us attribute any mIoU gap between them solely to \emph{prompt modality}: the SA-1B pipeline uses a fixed generic prompt while SACo-Gold uses concept text labels. A consistent
gap across all $\sigma$ levels reveals whether explicit text prompts provide more or less robust guidance to the visual encoder than the generic segmentation mode.

%--------------------------------------------------------------------
\subsection{Visual Degradation: Spatial Flips}
\label{sec:flips}
We test three modes: horizontal~(\textbf{h}), vertical~(\textbf{v}), and combined~(\textbf{both}). All prompts are held fixed. This degradation is applied to \emph{both} pipelines.


  The primary metric is \textbf{flip-back mIoU}: the predicted mask
  $\hat{\mathbf{M}}_\phi$ is mapped back to the original coordinate frame via
  the inverse transformation $T_\phi^{-1}$ before comparison against the clean
  baseline $\hat{\mathbf{M}}_0$:
  %
  \begin{equation}
    \mathrm{mIoU}_{\mathrm{fb}}(\phi)
    = \frac{1}{|\mathcal{D}|}
      \sum_{i \in \mathcal{D}}
      \mathrm{IoU}(T_\phi^{-1}(\hat{\mathbf{M}}_\phi^{(i)}),\;\hat{\mathbf{M}}_0^{(i)}).
    \label{eq:flipback}
  \end{equation}
  %
  A high $\mathrm{mIoU}_{\mathrm{fb}}$ indicates the model correctly segmented the object, but in the flipped coordinate frame equivariant behaviour. A low value indicates true semantic failure independent of coordinate alignment.

  To contextualise this, we also record the \textbf{direct mIoU}: the
  uncorrected comparison of $\hat{\mathbf{M}}_\phi$ against $\hat{\mathbf{M}}_0$
  in the original frame (Eq.~\ref{eq:miou-dataset}). Direct mIoU is expected to be low even for a perfectly equivariant model, since the mask will land in the flipped coordinate space. It serves only as the naive baseline from which the
  misregistration fraction is derived:
  %
  \begin{equation}
    \rho_\phi
    = \frac{
        \mathrm{mIoU}_{\mathrm{fb}}(\phi) - \mathrm{mIoU}_{\mathrm{direct}}(\phi)
      }{
        1 - \mathrm{mIoU}_{\mathrm{direct}}(\phi)
      },
    \label{eq:misreg}
  \end{equation}
  %
  where $\rho_\phi \approx 1$ indicates that nearly all apparent failure is spatial misregistration rather than semantic loss, and $\rho_\phi \approx 0$ indicates pure semantic failure.
%--------------------------------------------------------------------
\subsection{Evaluation Metrics}

Let $f_\theta$ denote SAM~3 and $(\mathbf{x}, \mathbf{p})$ an image prompt
pair.  Define the \emph{clean prediction}
$\hat{\mathbf{M}}_0 = f_\theta(\mathbf{x}, \mathbf{p})$ and the \emph{degraded prediction}
$\hat{\mathbf{M}}_\ell = f_\theta(\mathbf{x}', \mathbf{p}')$ at noise level $\ell$.  Each prediction is a set of instance masks $\{m_k\}$.

% Mean IoU (mIoU).

For each predicted mask $m_k \in \hat{\mathbf{M}}_\ell$ we compute the maximum
overlap with any mask in the clean baseline:
%
\begin{equation}
  \mathrm{IoU}(\hat{\mathbf{M}}_\ell,\hat{\mathbf{M}}_0)
  = \frac{1}{|\hat{\mathbf{M}}_\ell|}
    \sum_{m_k \in \hat{\mathbf{M}}_\ell}
    \max_{m_j \in \hat{\mathbf{M}}_0} \mathrm{IoU}(m_k, m_j).
  \label{eq:miou-instance}
\end{equation}
%
Dataset-level mIoU is the mean over all valid pairs $\mathcal{D}$:
%
\begin{equation}
  \mathrm{mIoU}(\ell)
  = \frac{1}{|\mathcal{D}|}
    \sum_{i \in \mathcal{D}}
    \mathrm{IoU}(\hat{\mathbf{M}}_\ell^{(i)}, \hat{\mathbf{M}}_0^{(i)}).
  \label{eq:miou-dataset}
\end{equation}
%
By construction $\mathrm{mIoU}(0) = 1.0$ for every condition.

%Sensitivity Index (SI).

To compare fragility across degradation types and prompt modalities on a common scale, we normalise the stability drop by the noise magnitude:
  %
  \begin{equation}
    S_\ell = \frac{\mathrm{mIoU}(\ell) - 1}{\ell},
    \label{eq:si}  \end{equation}
  %
  where $\ell$ is the noise level (number of character substitutions for linguistic noise; $\sigma$ for Gaussian noise). $S_\ell \leq 0$ always; a larger magnitude indicates greater fragility per unit of noise. Spatial flips are excluded from SI as their modes (h, v, both) are categorical and carry no
  meaningful numeric magnitude; flip results are characterised instead through direct mIoU, flip-back mIoU, and the misregistration fraction $\rho$
  (\S\ref{sec:flips}).
%-----------------------------------------------
\section{Experimental Setup}
\label{sec:expSetup}
We evaluate SAM~3~\cite{peng2025sam3} loaded from the official pretrained checkpoint via the HuggingFace Transformers library (\texttt{Sam3Model},
\texttt{Sam3Processor}). All inference is run in half-precision (\texttt{float16}). Instance segmentation outputs are post-processed with a confidence threshold of 0.5 and a mask binarisation threshold of 0.5,
following the default SAM~3 evaluation configuration.


Experiments were conducted across two platforms: a Kaggle notebook with dual NVIDIA Tesla T4 GPUs (16\,GB VRAM each), and a personal workstation with a
single NVIDIA RTX 2060 GPU (6\,GB VRAM).


For each image prompt pair, a single forward pass is performed and the resulting instance masks are encoded in COCO RLE format and appended to a JSONL
file, one file per degradation level. This design allows each condition to be reproduced independently. All mIoU, SI, and $\rho$ values are computed offline
by comparing each degraded-condition JSONL file against the corresponding clean-condition JSONL file.


Datasets are described in detail in \S\ref{sec:datasets}. The effective
evaluation counts after baseline matching are 9,442--9,443 image pairs per degradation level for the SA-1B pipeline, and 8,046--8,048 pairs for the SACo-Gold pipeline.


All experiments are run with a single random seed as gaussian noise sampling and character substitutions involve stochastic operations. 

%---------------------------------------------------------------------
\section{Results}
\label{sec:results}

\subsection{Gaussian Noise}
\label{sec:results-gaussian}

Table~\ref{tab:gaussian} shows output stability under Gaussian noise for both
pipelines.  The SA-1B pipeline loses stability steadily, dropping from mIoU\,=\,0.758 at $\sigma$\,=\,10 to
0.588 at $\sigma$\,=\,50.  The SACo-Gold pipeline
is notably more robust across all levels, reaching only 0.813 at the
highest severity.  Concept prompts appear to provide stronger guidance to the model's visual encoder when image quality degrades.

\begin{table}[h]
\centering
\small
\begin{tabular}{lcccc}
\hline
& \multicolumn{2}{c}{SA-1B} & \multicolumn{2}{c}{SACo-Gold} \\
$\sigma$ & mIoU & SI & mIoU & SI \\
\hline
10 & 0.758 & $-$0.0242 & 0.923 & $-$0.0077 \\
25 & 0.661 & $-$0.0136 & 0.869 & $-$0.0053 \\
50 & 0.588 & $-$0.0082 & 0.813 & $-$0.0037 \\
\hline
\end{tabular}
\vspace{3pt}
\caption{Output stability under Gaussian noise.
         mIoU measures consistency with the clean baseline;
         SI (Eq.~\ref{eq:si}) normalises the drop by $\sigma$.
         SACo-Gold is more robust than SA-1B at every level.}
\label{tab:gaussian}
\end{table}

%--------------------------------------------------------------------
\subsection{Text Degradation}
\label{sec:results-text}

Table~\ref{tab:text} shows output stability under linguistic noise on the
SACo-Gold pipeline.  A single character substitution ($\ell$\,=\,1) reduces
mIoU from 1.0 to 0.323 — a drop of 0.677, far larger than anything observed
under visual noise.  Two substitutions push mIoU below 0.15; three leave
almost no consistent output (mIoU\,=\,0.074).  SAM~3's concept-prompt
pipeline is highly sensitive to even minimal typographic corruption.

\begin{table}[h]
\centering
\small
\begin{tabular}{lcc}
\hline
$\ell$ & mIoU & SI \\
\hline
1 & 0.323 & $-$0.677 \\
2 & 0.148 & $-$0.426 \\
3 & 0.074 & $-$0.309 \\
\hline
\end{tabular}
\vspace{3pt}
\caption{Output stability under linguistic noise (SACo-Gold pipeline).
         mIoU\,=\,1.0 at $\ell$\,=\,0 by construction.
         A single character substitution causes a 68\% drop in SI.}
\label{tab:text}
\end{table}

%--------------------------------------------------------------------

\subsection{Spatial Flips}
\label{sec:results-flips}

Table~\ref{tab:flips} reports direct mIoU, flip-back mIoU, and the
misregistration fraction $\rho$ for each flip mode and pipeline.  Direct mIoU
is low in all cases, as expected for an unaligned comparison.  Flip-back mIoU
recovers much of that gap, indicating that the model successfully segments the
target object in the flipped frame.  For the SA-1B pipeline, $\rho$ ranges
from 0.407 (vertical) to 0.675 (horizontal), meaning 41\%--68\% of the
apparent failure is purely geometric misregistration.  The SACo-Gold pipeline
is more spatially equivariant across all modes: $\rho$ reaches 0.865 for
horizontal flips, indicating that 86\% of the apparent failure is
misregistration rather than semantic loss.  Horizontal flips are consistently
easier to recover than vertical or combined flips in both pipelines, likely
because horizontal symmetry is more common in natural images.

\begin{table}[h]
\centering
\small
\begin{tabular}{llccc}
\hline
Pipeline & Mode & Direct & Flip-back & $\rho$ \\
\hline
\multirow{SA-1B}
  & Horizontal & 0.094 & 0.705 & 0.675 \\
  & Vertical   & 0.055 & 0.439 & 0.407 \\
  & Both       & 0.031 & 0.431 & 0.413 \\
\hline
\multirow{SACo-Gold}
  & Horizontal & 0.204 & 0.892 & 0.865 \\
  & Vertical   & 0.134 & 0.692 & 0.644 \\
  & Both       & 0.090 & 0.685 & 0.654 \\
\hline
\end{tabular}
\vspace{3pt}
\caption{Spatial flip results.  Direct mIoU compares flipped predictions
         against the clean baseline without alignment.  Flip-back mIoU
         applies the inverse transform before comparison (Eq.~\ref{eq:flipback}).
         $\rho$ (Eq.~\ref{eq:misreg}) measures what fraction of the
         apparent failure is geometric misregistration rather than
         semantic loss.}
\label{tab:flips}
\end{table}

%--------------------------------------------------------------------
\subsection{Linguistic Fragility by Prompt Length}
\label{sec:results-text-length}

Table~\ref{tab:prompt-length-si} breaks down text degradation by prompt
word count.  Single-word concepts are the most fragile: one substitution
reduces mIoU to just 0.027 (SI\,=\,$-$0.973), nearly complete collapse.
Two-word prompts are more resilient but still degrade sharply.
Three-or-more-word phrases retain substantially more stability at $\ell$\,=\,1
(mIoU\,=\,0.483), consistent with their lower effective corruption rate
of 6\% compared to 17\% for single-word concepts.  The result confirms
that fragility under character-level noise scales with the fraction of the prompt that is corrupted, not just the absolute number of substitutions.

\begin{table}[h]
\centering
\small
\begin{tabular}{lcccccc}
\hline
& \multicolumn{2}{c}{$\ell$=1} & \multicolumn{2}{c}{$\ell$=2} & \multicolumn{2}{c}{$\ell$=3} \\
B & mIoU & SI & mIoU & SI & mIoU & SI \\
\hline
1   & 0.027 & $-$0.973 & 0.009 & $-$0.496 & 0.002 & $-$0.333 \\
2   & 0.208 & $-$0.792 & 0.053 & $-$0.474 & 0.017 & $-$0.328 \\
3+  & 0.483 & $-$0.517 & 0.253 & $-$0.373 & 0.135 & $-$0.288 \\
\hline
\end{tabular}
\vspace{3pt}
\caption{mIoU and SI per prompt-length bucket under linguistic degradation (SACo-Gold pipeline).  Effective corruption rates at $\ell$\,=\,1 are 17\%, 11\%, and 6\% for 1-, 2-, and
         3+-word buckets respectively.}
\label{tab:prompt-length-si}
\end{table}
%--------------------------------------------------------------------

\subsection{Visual vs.\ Linguistic Fragility}
\label{sec:results-crossmodal}

Table~\ref{tab:crossmodal} pairs Gaussian and text noise at three
matched severity tiers.
At every tier, a single-digit number of character substitutions does more
damage to output consistency than the strongest Gaussian corruption tested.
At mild severity, text degrades mIoU to 0.323
while Gaussian keeps both pipelines above 0.758. At severe severity, text reaches 0.074 while Gaussian stays above 0.588.
\begin{table}[h]
\centering
\footnotesize
\setlength{\tabcolsep}{5pt}
\begin{tabular}{lccc}
\hline
Tier & SA-1B\textsuperscript{G} & SACo\textsuperscript{G} & SACo\textsuperscript{T} \\
\hline
Mild     & 0.758 & 0.923 & 0.323 \\
Moderate & 0.661 & 0.869 & 0.148 \\
Severe   & 0.588 & 0.813 & 0.074 \\
\hline
\end{tabular}
\vspace{3pt}
\caption{mIoU at matched severity tiers.
         G\,=\,Gaussian ($\sigma$\,=\,10/25/50),
         T\,=\,text ($\ell$\,=\,1/2/3).
         Text corruption (right column) degrades stability far below
         both Gaussian conditions at every tier.}
\label{tab:crossmodal}
\end{table}

%-------------------------------------------------------------------- 
\section{Conclusion}
\label{sec:conc}
%--------------------------------------------------------------------
{\small
\bibliographystyle{ieee}
\bibliography{egbib}
}
%---------------------------------------------------------------------

\section{Appendix}

\subsection{Geometric Degradation (Deferred)}
\label{sec:geo-deferred}

Bounding-box prompt jitter(expanding or shifting ground-truth boxes by
10\%--20\%)was implemented but excluded from the main evaluation.  Prior work already characterises SAM~1/2 geometric sensitivity
thoroughly~\cite{moskalenko2026breps}; our primary contribution is the
text-prompt and cross-modal axes, which are entirely unstudied on SAM~3. A geometric prong for SAM~3, including combined text+geometric noise, is left to future work.

\end{document}
